% Präambel
\documentclass[%
12pt,							% Schriftgröße
a4paper,						% Papierformat
oneside, 						% einseitiges (oneside) oder zweiseitiges (twoside) Dokument
listof=totoc, 					% Tabellen- und Abbildungsverzeichnis ins Inhaltsverzeichnis
bibliography=totoc,				% Literaturverzeichnis ins Inhaltsverzeichnis aufnehmen
titlepage, 						% Titlepage-Umgebung statt \maketitle
headsepline, 					% horizontale Linie unter Kolumnentitel
%abstracton,					% Überschrift beim Abstract einschalten, Abstract muss dazu in {abstract}-Umgebung stehen
DIV=12,							% Satzspiegeleinstellung, 12 ist Standar bei KOMA
%BCOR=6mm,						% Bindekorrektur, die den Seitenspiegel um 6mm nach rechts verschiebt,
cleardoublepage=empty,			% Stil einer leeren eingefügten Seite bei Kapitelwechsel
parskip,							% Absatzabstand bei Absatzwechsel einfügen
ngerman
]{scrbook}		
\usepackage{scrhack}
\usepackage[utf8]{inputenc} 	% ermöglicht die direkte Eingabe von Umlauten
\usepackage[T1]{fontenc} 		% Ausgabe aller zeichen in einer T1-Codierung (wichtig für die Ausgabe von Umlauten!)
\usepackage{babel} 	% deutsche Trennungsregeln und Übersetzung der festcodierten Überschriften
\setlength{\parindent}{0ex} 	% bei neuem Abschnitt nicht einrücken
%------
% Folgende Einstellungen entsprechen den Vorgaben der Leitlinien
\usepackage[onehalfspacing]{setspace}
% Ende Leitlinien
%------
%------
% Folgende Einstellungen sind bei größeren Arbeiten mit viel Text zu empfehlen.
% Hierbei oben DIV=16 einstellen und Zeile \usepackage[onehalfspacing]{setspace} auskommentieren.
%\linespread{1.2}\selectfont     % Zeilenabstand erhöhen - größere Werte als 1.2 nicht verwenden!!
% Ende Einstellung große Arbeiten mit viel Text.
%------

\usepackage{siunitx}			% Vereinfachte Eingabe von Einheiten in Formeln
\sisetup{
	number-unit-product = \;,
	inter-unit-product = \:,
	exponent-product = \cdot,
	output-decimal-marker = {,}
}

\usepackage{graphicx}  			% Einbinden von Grafiken erlauben
\usepackage[format=hang,		% Formatierungen von Unter- / Überschriften
font=normal,
labelfont=bf,
justification=RaggedRight,
singlelinecheck=true,
aboveskip=1mm
]{caption}

\usepackage[backend=biber, %% Hilfsprogramm "biber" beim Compilieren nutzen (statt "biblatex" oder "bibtex")
style=alphabetic, %% Zitierstil (siehe Dokumentation)
natbib=true, %% Bereitstellen von natbib-kompatiblen Zitierkommandos
hyperref=true, %% hyperref-Paket verwenden, um Links zu erstellen
]{biblatex}
\addbibresource{literature/literatur1.bib} %% Einbinden der bib-Datei. Endung .bib unbedingt ergänzen
\addbibresource{literature/literatur2.bib} %% Einbinden mehrerer bib-Dateien mit zusätzlichem \addbibresource - Befehl

% Folgende Zeilen sind auszukommentieren, falls runde Klammern und ein vgl. bei Zitaten erscheinen sollen.
%\makeatletter
%\renewcommand{\@cite}[2]{(vgl. {#1\if@tempswa , #2\fi})} 
%\renewcommand{\@biblabel}[1]{(#1)}
%\makeatother

\usepackage{pdfpages}

\usepackage{enumitem}			% Erlaubt Änderung der Nummerierung in der Umgebung enumerate

\usepackage{amsmath}			% Ergänzungen für Formeln
\usepackage{textcomp} 			% zum Einsatz von Eurozeichen u. a. Symbolen
\usepackage{eurosym}			% bessere Darstellung Euro-Symbol mit \euro
\newcommand*\diff{\mathop{}\!\mathrm{d}}	% Differentialzeichen
\newcommand*\Diff[1]{\mathop{}\!\mathrm{d^#1}} % Differentialzeichen höherer Ableitung
\newcommand*\jj{\mathop{}\!\mathrm{j}}	% Komplexe Zahl j
\usepackage{listings}

\usepackage[					% Einstellunge Paket hyperref
hyperfootnotes=false,			% im pfd-Output Fußnoten nicht verlinken
hidelinks
]{hyperref}

\usepackage{makeidx}			% Paket zur Erstellung eines Index
\usepackage[intoc]{nomencl} 	% zur Erstellung des Abkürzungsberzeichnisses+
\usepackage{underscore}

\usepackage[					% Einstellungen für Fußnoten
bottom,							% Ausrichtung unten
multiple,						% Trennung durch Seperator bei mehreren Fußnoten
hang,
marginal
]{footmisc}

\usepackage{calc}				% Paket zum Berechnen von Längen z.B. 0.8\linewidth

\usepackage{xcolor} 			% einfache Verwendung von Farben in nahezu allen Farbmodellen

\usepackage{listings}			% Darstellung von Quellcode mit den Umgebungen {lstlisting}, \lstinline und \lstinputlisting
\lstset{literate=				% Damit können Umlaute innerhalb Listings geschrieben werden
	{Ö}{{\"O}}1
	{Ä}{{\"A}}1
	{Ü}{{\"U}}1
	{ß}{{\ss}}1
	{ü}{{\"u}}1
	{ä}{{\"a}}1
	{ö}{{\"o}}1
}
\definecolor{mygreen}{rgb}{0,0.6,0}
\definecolor{mygray}{rgb}{0.5,0.5,0.5}
\definecolor{mymauve}{rgb}{0.58,0,0.82}
\lstset{ %
	backgroundcolor=\color{white},   % choose the background color; you must add \usepackage{color} or \usepackage{xcolor}; should come as last argument
	basicstyle=\footnotesize,        % the size of the fonts that are used for the code
	breakatwhitespace=false,         % sets if automatic breaks should only happen at whitespace
	breaklines=true,                 % sets automatic line breaking
	captionpos=t,                    % sets the caption-position to (b) bottom or (t) top
	commentstyle=\color{mygreen},    % comment style
	deletekeywords={...},            % if you want to delete keywords from the given language
	escapeinside={\%*}{*)},          % if you want to add LaTeX within your code
	escapeinside={(*@}{@*)},
	extendedchars=true,              % lets you use non-ASCII characters; for 8-bits encodings only, does not work with UTF-8
	frame=none,	                   	% "single" adds a frame around the code; "none"
	keepspaces=true,                 % keeps spaces in text, useful for keeping indentation of code (possibly needs columns=flexible)
	keywordstyle=\color{blue},       % keyword style
	language=[LaTeX]TeX,             % the language of the code
	morekeywords={*,nomenclature},   % if you want to add more keywords to the set
	numbers=left,                    % where to put the line-numbers; possible values are (none, left, right)
	numbersep=5pt,                   % how far the line-numbers are from the code
	numberstyle=\tiny\color{mygray}, % the style that is used for the line-numbers
	rulecolor=\color{black},         % if not set, the frame-color may be changed on line-breaks within not-black text (e.g. comments (green here))
	showspaces=false,                % show spaces everywhere adding particular underscores; it overrides 'showstringspaces'
	showstringspaces=false,          % underline spaces within strings only
	showtabs=false,                  % show tabs within strings adding particular underscores
	stepnumber=1,                    % the step between two line-numbers. If it's 1, each line will be numbered
	stringstyle=\color{mymauve},     % string literal style
	tabsize=2,	                   % sets default tabsize to 2 spaces
	title=\lstname                   % show the filename of files included with \lstinputlisting; also try caption instead of title
}

\makeindex						% Indexverzeichnis erstellen
\makenomenclature				% Abkürzungsverzeichnis erstellen

% -----------------------------------------------------------------------------------------------------------------
% Zum Aktualisieren des Abkürzungsverzeichnisses (Nomenklatur) bitte auf der Kommandozeile folgenden Befehl aufrufen :
% makeindex <Dateiname>.nlo -s nomencl.ist -o <Dateiname>.nls
% Oder besser: Kann in TexStudio unter Tools-Benutzer als Shortlink angelegt werden
% Konfiguration unter: Optionen-Erzeugen-Benutzerbefehle: makeindex -s nomencl.ist -t %.nlg -o %.nls %.nlo
% -----------------------------------------------------------------------------------------------------------------

% Hier die persönlichen Daten eingeben:

\newcommand{\titel}{Programmentwurf}
\newcommand{\untertitel}{Entwicklung eines rudimentären Zeitscheibensystems in C}
\newcommand{\arbeit}{}
\newcommand{\studiengang}{Embedded Systems }
\newcommand{\studienrichtung}{Automotive Engineering}
\newcommand{\studienschwerpunkt}{}
\newcommand{\autor}{Fabian Klotz}
\newcommand{\matrikelnr}{}
\newcommand{\kurs}{TSA22}
\newcommand{\firma}{}
\newcommand{\abgabe}{09.04.2023}
\newcommand{\betreuerdhbw}{Prof. Dr. Florian Leitner-Fischer}
\newcommand{\betreuerfirma}{}
\newcommand{\jahr}{2023}			% für Angabe im Copyright-Vermerk der Titelseite

% Folgende Zeilen definieren Abkürzungen, um Befehle schneller eingeben zu können
\newcommand{\ua}{\mbox{u.\,a.\ }}
\newcommand{\zB}{\mbox{z.\,B.\ }}
\newcommand{\bs}{$\backslash$}

% Folgende Zeilen weden benötigt, um Tikz und PGF-Plot-Grafiken einzubinden
\usepackage{pgfplots}
\usepackage{pgfplotstable}
\pgfplotsset{compat=newest,width=0.6\linewidth}
\usepgfplotslibrary{smithchart}
\usepackage{tikz}						% Tikz sollte nach Listings Pakete geladen werden.
\usetikzlibrary{arrows}

\hyphenation{Schrift-ar-ten}


% -------------------------------------------------------------------------------------------
%                     Beginn des Dokumenteninhalts
% -------------------------------------------------------------------------------------------
\begin{document}
\let\texteuro\euro						% Eingabe \texteuro, € oder \euro erzeugt gleiches Ergebnis
\setcounter{secnumdepth}{3}				% Nummerierungstiefe fürs Inhaltsverzeichnis
\setcounter{tocdepth}{3}
\sffamily								% für die Titelei serifenlose Schrift verwenden

% ------------------------------ Titelei -----------------------------------------------------

\thispagestyle{plain}
\hypersetup{pageanchor=false}
\begin{titlepage}
\enlargethispage{4.0cm}
\sffamily 								% Serifenlose Grundschrift für die Titelseite einstellen

\parbox{0.5\linewidth}{
\begin{flushleft}
% Hier ggf. ein Logo der Firma
\end{flushleft}
}
\parbox{0.5\linewidth}{
\begin{flushright}
	\includegraphics[width=0.4\linewidth]{images/DHBW_d_R_FN_46mm_4c}\\[5ex]
\end{flushright}
}
				

\begin{center}

{\fontsize{20.74pt}{24pt}\selectfont
\textbf{\titel}\\[1.5ex]}
{\fontsize{14pt}{17pt}\selectfont
\textbf{\untertitel}\\[5ex]}
{\fontsize{17pt}{20pt}\selectfont
\textbf{\arbeit}\\[2ex]}
{\fontsize{14pt}{17pt}\selectfont
Studiengang: \studiengang\\[2ex]}
{\fontsize{12pt}{14pt}\selectfont
Studienrichtung: \studienrichtung\\[1ex]
Duale Hochschule Baden-Württemberg Ravensburg, Campus Friedrichshafen\\[5ex]
von\\[1ex]
\autor\\[15ex]}


\end{center}

\begin{flushleft}
{\fontsize{12pt}{14pt}\selectfont
\begin{tabular}{ll}
Abgabedatum:					& \quad \abgabe \\
Kurs: 							& \quad \kurs \\
Studiengangsleiter: & \quad \betreuerdhbw \\ [2ex]
\end{tabular}
}
\end{flushleft}
%%%%% Nachfolgende Zeilen einkommentieren, wenn Copyrightvermerk gewünscht ist
%\begin{flushleft}
%{\fontsize{11pt}{13pt}\selectfont
%Copyrightvermerk:\\
%Dieses Werk einschließlich seiner Teile ist \textbf{urheberrechtlich geschützt}. Jede Verwertung außerhalb der engen Grenzen des Urheberrechtgesetzes ist ohne Zustimmung des Autors unzulässig und strafbar. Das gilt insbesondere für Vervielfältigungen, Übersetzungen, Mikroverfilmungen sowie die Einspeicherung und Verarbeitung in elektronischen Systemen.
%}
%\end{flushleft}
%\begin{flushright}
%{\fontsize{11pt}{13pt}\selectfont \copyright{} \jahr }
%\end{flushright}
\end{titlepage}

\cleardoublepage
\hypersetup{pageanchor=true}
 				% erzeugt die Titelseite
\pagenumbering{roman}					% kleine, römische Seitenzahlen für Titelei
%% Ggf. folgende Zeile auskommentieren, falls der Sperrvermerk gewünscht ist.
%\chapter*{Sperrvermerk} %*-Variante sorgt dafür, das der Sperrvermerk nicht im Inhaltsverzeichnis auftaucht
%gemäß Ziffer 1.1.13 der Anlage 1 zu §§ 3, 4 und 5  der Studien- und Prüfungsordnung für die Bachelorstudiengänge im Studienbereich Technik der Dualen Hochschule Baden-Württemberg vom 29.09.2017 in der Fassung vom 25.07.2018:
%
%Der Inhalt dieser Arbeit darf weder als Ganzes noch in Auszügen Personen außerhalb des Prüfungsprozesses und des Evaluationsverfahrens zugänglich gemacht werden, sofern keine anders lautende Genehmigung vom Dualen Partner vorliegt.
%
%Musterstadt, den \today \\[4ex]
%
%\rule[-0.2cm]{5cm}{0.5pt} \\
%
%\textsc{\autor} \\[10ex]

\chapter*{Erklärung} %*-Variante sorgt dafür, dass die Erklärung nicht im Inhaltsverzeichnis auftaucht

gemäß Ziffer 1.1.13 der Anlage 1 zu §§ 3, 4 und 5  der Studien- und Prüfungsordnung für die Bachelorstudiengänge im Studienbereich Technik der Dualen Hochschule Baden-Württemberg vom 29.09.2017 in der Fassung vom 25.07.2018.

Ich versichere hiermit, dass ich meinen Programmentwurf mit dem Thema:

\begin{quote}
	\textit{ \untertitel }
\end{quote}

selbstständig verfasst und keine anderen als die angegebenen Quellen und Hilfsmittel benutzt habe.\\[6ex]

Weingarten, den \today \\[1ex]

\rule[-0.2cm]{5cm}{0.5pt} \\

\autor \\[10ex]

\rmfamily

\thispagestyle{empty}

\cleardoublepage

 				% Einbinden der eidestattlichen Erklärung
\chapter*{Kurzfassung} %*-Variante sorgt dafür, das Abstract nicht im Inhaltsverzeichnis auftaucht

Im Rahmen des Programmentwurfs der Vorlesung Programmieren soll ein rudimentäres Zeitscheibensystem entworfen werden.
Hierbei sind einige Anforderungen zu erfüllen, wie zum Beispiel eine sehr große Variabilität des Programms, sodass mit möglichst wenig Aufwand neue Tasks angelegt, oder zu anderen Zykluszeiten umgesetzt werden können. 

Das vorliegende Programm basiert hierbei auf dem Konzept des Threadings um die zeitliche Abfolge zu organisieren. 
In den einzelnen Tasks können Funktionen mit einer bestimmten Priorität, also Vorgabe der auszuführenden Reihenfolge, konfiguriert werden.
Außerdem können sich die einzelnen Funktionen selbst für die Ausführung in einem Task mit einer gewissen Priorität registrieren.
Braucht eine Funktion länger für die Ausführung wie die angegebene Zykluszeit, wird der Task erst dann neu gestartet, wenn er regulär wieder an der Reihe ist.

Abschließend wird der Worst-Case Jitter der einzelnen Tasks berechnet.

\cleardoublepage
   			% Einbinden des Abstracts

\tableofcontents						% Erzeugen des Inhalsverzeichnisses
\cleardoublepage

% --------------------------------------------------------------------------------------------
%                    Inhalt der Bachelorarbeit
%---------------------------------------------------------------------------------------------
\pagenumbering{arabic}					% arabische Seitenzahlen für den Hauptteil

\rmfamily

\include{chapter/einleitung}
\include{chapter/grundlagen}
\chapter{Konzeptentwurf}
\label{cha:Konzeptentwurf}

Um ein passendes Konzept für die Umsetzung des Programms zu finden, wurden im Voraus einige Überlegungen, vor allem zur Steuerung der zeitlichen Abfolge angestellt.
Mit dem Konzept des Threadings kann sowohl die zeitliche Abfolge, als auch der Aufruf der einzelnen Funktionen in den Tasks effizient gesteuert werden. 

Die Abbildung  soll Threading bildhaft darstellen\footnote[1]{https://www.baeldung.com/wp-content/uploads/sites/4/2020/07/multithreading.png}	.

\begin{figure}[hbt]				% here, bottom, top
	\centering						% Zentrierung
	\includegraphics[width=0.6\linewidth]{images/multithreading}	
	\label{fig:multithreading}
	\caption{Threading}
	{\raggedright
	 }
	
\end{figure}

Der \glqq Main Thread\grqq \space startet mehrere untergeordnete Threads. Alle Threads greifen auf die gleichen Ressourcen und globalen Variablen zu. 

Am Anfang des Programms wird definiert, wie viele Tasks (bzw. Threads) es geben soll. Die genannte Anzahl Threads wird draufhin gestartet. Jeder Thread bekommt eine eigene Verkettete Liste, in der die einzelnen ausführbaren Funktionen mit einer Priorität gespeichert werden. Wird eine Funktion zur Liste hinzugefügt, sortiert ein Algorithmus die Funktionen nach Priorität. 
Die Threads durchlaufen alle die gleiche Funktion, in der durch die Listen iteriert wird und die Funktionen abgearbeitet werden. Anschließend wird der Thread so lange schlafen gelegt, bis der nächste reguläre Zyklus ansteht. 
Durch die Verwendung von Mutex werden Race-Conditions verhindert.
Die Threads laufen jeweils in einer Endlosschleife, damit das Programm erste terminiert wird, wenn es der User abbricht. 
Wird das Programm von außen terminiert, wird der von den Listen allozierte Speicher wieder freigegeben. 
\chapter{Umsetzung}
\label{cha:Umsetzung}
\section{pthread}
Um erste Versuche mit Threading zu starten, wird ein Programm erstellt, das zwei Threads startet. 
Dabei wurde die Bibliothek 'pthread.h' verwendet. 

\texttt{pthread} steht für ' POSIX Threads' und ist eine API (Application Programming Interface), die es ermöglicht, in C und C++ Threads zu erstellen und zu verwalten. POSIX steht dabei für "Portable Operating System Interface", was bedeutet, dass diese API auf vielen verschiedenen Betriebssystemen verfügbar ist und somit portabel ist. Pthreads werden unter UNIX-basierten Betriebssystemen wie MacOS oder Linux unterstützt. 

Die \texttt{pthread}-API stellt eine Reihe von Funktionen bereit, mit denen Threads erstellt, gesteuert und synchronisiert werden können. Einige wichtige Funktionen der \texttt{pthread}-API sind:

\begin{itemize}
\item \verb|pthread_create|: Diese Funktion erstellt einen neuen Thread und führt eine bestimmte Funktion aus, die dem Thread als Argument übergeben wird.
\item \verb|pthread_join|: Diese Funktion blockiert den aufrufenden Thread, bis der angegebene Thread beendet wurde.
\end{itemize}

Die \texttt{pthread}-API ist sehr leistungsfähig und wird häufig in Anwendungen eingesetzt, die viele gleichzeitig ausgeführte Threads erfordern, wie beispielsweise Serveranwendungen oder Multimedia-Anwendungen\footnote[2]{https://www.cs.cmu.edu/afs/cs/academic/class/15492-f07/www/pthreads.html}. 


Im fertigen Programm wird automatisch die richtige Anzahl an Threads erstellt. 
Die untenstehende for-Schleife wird so oft wiederholt, wie es Tasks im Programm geben soll.
In der Schleife wird jeweils 200ms gewartet, damit die Threads leicht zeitversetzt laufen, um später die Ausgabe auf der Konsole lesbarer zu machen.

\begin{lstlisting}[language=C]
pthread_t threads[num_threads];
int thread_args[num_threads];

for (int i = 0; i < num_threads; i++) {
	thread_args[i] = i;
	pthread_create(&threads[i], NULL, thread_function, &thread_args[i]); 
	wait_200ms();
}
\end{lstlisting}

Beim erstellen eines Threads wird immer die \texttt{thread_function} aufgerufen. Hier wird das Grundgerüst der einzelnen Threads festgelegt. Es wird eine \texttt{while(1)}, also eine Endlosschleife angelegt. In dieser Schleife wird zuerst ein Timer gestartet und danach mit einer weiteren Funktion überprüft, ob im Konfigurationsarray eine Funktion für diesen Task hinterlegt ist. Sollte das der Fall sein, werden die Aufgaben je nach Priorität abgearbeitet. Laufen die auszuführenden, konfigurierten Funktionen zu lange, also länger wie die vorgegebene Zykluszeit, wird so lange gewartet, bis der Task wieder regulär an der Reihe ist. 

\section{Konfiguration der Funktionen}
\subsection{Konfigurationsarray}

Das Konfigurationsarray um die Funktionen den Tasks zuzuordnen ist folgendermaßen aufgebaut:
\begin{lstlisting}[language=C]
double array[NUM_FUNC][NUM_TASKS] = {
    //          T1  T2  T3  T4
    /*add*/    {3,  0,  0,  0},
    /*multi*/  {2,  2,  0,  0},
    /*div*/    {1,  1,  1,  0},
    /*mod*/    {0,  0,  0,  1},
    /*div2*/   {0,  3,  0,  0}
};
\end{lstlisting}

Wie an den Kommentaren zu sehen ist, geben die Zeilen immer die Funktionen an, welche in den Tasks auf gerufen werden sollen und die Spalten die dazugehörigen Tasks. Die Zahlen stehen hierbei für die Priorität, also die Reihenfolge. (1 = höchste Priorität, 0 = nicht Konfiguriert) 

\subsection{Selbstregistrierung von Funktionen zur Ausführung in einem Task}
Damit wie gefordert eine Selbstregistrierung der Funktionen zur Ausführung in einem Task (bzw. Thread) ausgeführt werden kann, wird beim Aufruf der mathematischen Funktionen unter anderem ein Übergabewert übermittelt, der dazu führt, dass über eine if-Bedingung eine Funktion \texttt{registration()} aufgerufen wird. Als Übergabewert erwartet diese Funktion den Task, die auszuführende Funktion und die Priorität. 
\begin{lstlisting}[language=C]
void registration(int task, int priority, int func){
        printf("\nTask %d registriert Funktion%d mit Priorität %d \n\n", task, func, priority);
        array[func-1][task-1] = priority;
}
\end{lstlisting}

\subsection{Ausführung der Funktionen in den Tasks}

Die Funktion \texttt{checkForTasks} läuft folgendermaßen ab:
\begin{lstlisting}[language=C]
void checkForTasks(int thread, double array[][NUM_TASKS]) {
    for(int j = 0; j < NUM_FUNC; j++){
        double *p = &array[0][thread-1];
        int i;
        for (i = 0; i < NUM_FUNC; i++) {
            if (*p == j+1) {
                fktPointer[i%NUM_FUNC](a, b, 0);
                printf("\n");
            }
            p = p + NUM_TASKS;
        }
    }
}
\end{lstlisting}

Beim Aufruf der Funktion wird der Thread und das Konfigurationsarray übergeben. Eine for-Schleife so oft durchlaufen, wie es Funktionen für die Tasks gibt. Danach wird ein Pointer auf das Array erstellt, der auf den ersten, für den Thread relevanten Eintrag zeigt. 
Es wird eine zweite for-Schleife erstellt, die ebenso so oft durchlaufen wird, wie es relevante Funktionen für die Tasks gibt. Danach wird überprüft, ob der Pointer auf ein Element zeigt, das der Schleifenvariablen der äußeren for-Schleife + 1 entspricht. Sollte das der Fall sein, wird ein Pointerarray aufgerufen, in dem Funktionspointer für die jeweiligen Funktionen hinterlegt sind. Durch die Rest-berechnung, innere Schleifenvariable  modulo  Anzahl der Funktionen wird die richtige Stelle im Array ermittelt und ausgeführt. 

\chapter{Bedienung}
\label{cha:Bedienung}

\section{Kompilierung}
Zur Kompilierung des vorliegenden Programms muss beachtet werden, dass pthreads nur auf UNIX-basierten Betriebssystemen verfügbar ist. 
Getestet wurde die Kompilierung ausschließlich in folgendem Docker Container:\newline
\texttt{leitnerfischerdhbw/es-ubuntu-x86}\newline\newline
Pthread muss manuell für das Kompilieren hinzugefügt werden. Der Befehl für die Kompilierung lautet als: \newline
\texttt{gcc -pthread  'Dateiname'}

\section{Änderung der Anzahl auszuführender Tasks}
Um die Anzahl der Tasks zu ändern muss zuerst folgende Zeile geändert werden und auf die gewünschte Anzahl erhöht oder erniedrigt werden:\newline\texttt{#define NUM_TASKS 4}\newline
Danach muss die Zykluszeit festgelegt werden. Hierfür wird in diesem Array ein Element mit der gewünschten Zeit hinzugefügt werden (Zeit in ms):
\newline\texttt{int cycleTimes[NUM_TASKS] = \{1000, 5000, 10000, 100000\};} \newline
Zusätzlich sollte im Konfigurationsarray eine Spalte mit den gewünschten Funktionen und Priorität zur Ausführung in diesem Task hinzugefügt werden. Sollte das Array nicht erweitert werden, kann es zu unschönen Ergebnissen führen, da die nicht initialisierten Elemente nicht automatisch '0' sind.  
\section{Änderung der Zykluszeiten der Tasks}
Um die Die Zykluszeiten zu ändern muss lediglich in diesem Array die gewünschte Zeit geändert werden (Zeit in ms).
\newline\texttt{int cycleTimes[NUM_TASKS] = \{1000, 5000, 10000, 100000\};} \newline

\section{Hinzufügen von Funktionen}
Um eine Funktion hinzuzufügen muss im ersten Schritt folgender Wert erhöht bzw. erniedrigt werden:
\newline\texttt{#define NUM_FUNC 5}
\newline
Anschließend wird die neue Funktion analog zu den bestehenden Funktionen angelegt. Die Übergabewerte sind hierbei die double Werte a und b und eine int Variable um eine potentielle Selbstregistrierung zu ermöglichen. Im Array \texttt{double (*fktPointer[NUM_FUNC])(double, double, double)} muss außerdem ein neues Element angelegt werden, das gleich zu benennen ist, wie die erstellte Funktion
Abschließend muss das Konifgurationsarray um eine (oder mehr) Zeilen erweitert und initialisiert werden. 

\chapter{Jitter}
\label{cha:Worst Case Jitter}

Obwohl Multithreading ein sehr Effizientes Konzept der Informatik ist, kann es sein, dass die Systemauslastung so hoch ist, dass die Threads nicht gleichzeitig ausgeführt werden können. Das kann verschiedene Gründe, wie zum Beispiel der Anzahl der verfügbaren Prozessorkerne, die Auslastung des Arbeitsspeichers etc., haben. Im schlimmsten Fall, sind so wenige Systemressourcen vorhanden, dass alle Threads nacheinander ausgeführt werden müssen. In diesem Fall würde sich der Jitter folgendermaßen berechnen: \newline\newline
\sum Zykluszeiten
\par
In diesem Fall ist der Worst-Case-Jitter also 135 Sekunden. 
Um die Behauptung zu validieren, dass der Jitter von der Verfügbarkeit von Systemressourcen abhängt, wurde während der Ausführung des Programms ein CPU-Benchmark gestartet. Tatsächlich kannte beobachtet werden, dass die Ausführungszeit der einzelnen Threads teilweise deutlich erhöht wurde. 
Getestet wurde auf einem MacOS System mit Apple M1 und 16gb RAM. Die maximal zu beobachtende Abweichung ist 99ms. 


%chapter{Verifikation und Diskussion}
\label{cha:Verifikation}

... Verifikation, Auswertung, Lösungsbewertung, Diskussion der Ergebnisse
%\chapter{Tests}
\label{cha:tests}

Um das vorliegende Programm zu testen wird dieses als erstes Kompiliert, um sicherzugehen, dass keine Kompilerwahrnungen oder Errors auftreten. 

% ---- Literaturverzeichnis ----------
\interlinepenalty 10000					% Verhindert einen Umbruch mitten in Literatureinträgen
\printbibliography						% Erstellen des Literaturverzeichnisses

% -----Ausgabe aller Verzeichnisse ---
\setlength{\parskip}{0.5\baselineskip}
\renewcommand{\indexname}{Sachwortverzeichnis}
\printindex								% Erzeugen des Indexverzeichnises
\addcontentsline{toc}{chapter}{\indexname}
% alle Abkürzungen, die in der Arbeit verwendet werden. Die Alphabetische Sortierung übernimmt Latex. Nachfolgend sind Beispiele genannt, welche nach Bedarf angepasst, gelöscht oder ergänzt werden können.

% Bei den unten stehenden Formelzeichen ist erläutert, wie explizite Sortierschlüssel über den Inhalt der eckigen Klammer angegeben werden.

% Allgemeine Abkürzungen %%%%%%%%%%%%%%%%%%%%%%%%%%%%
%\nomenclature{Abb.}{Abbildung}
%\nomenclature{bzw.}{beziehungsweise}
%\nomenclature{DHBW}{Duale Hochschule Baden-Württemberg}
%\nomenclature{ebd.}{ebenda}
%\nomenclature{et al.}{at alii}
%\nomenclature{etc.}{et cetera}
%\nomenclature{evtl.}{eventuell}
\nomenclature{f.}{folgende Seite}
\nomenclature{ff.}{fortfolgende Seiten}
%\nomenclature{ggf.}{gegebenenfalls}
%\nomenclature{Hrsg.}{Herausgeber}
%\nomenclature{Tab.}{Tabelle}
%\nomenclature{u. a.}{unter anderem}
%\nomenclature{usw.}{und so weiter}
\nomenclature{vgl.}{vergleiche}
%\nomenclature{z. B.}{zum Beispiel}
%\nomenclature{z. T.}{zum Teil}

% Dateiendungen %%%%%%%%%%%%%%%%%%%%%%%%%%%%%%%%%%%%
\nomenclature{EMF}{Enhanced Metafile}
\nomenclature{JPG}{Joint Photographic Experts Group}
\nomenclature{PDF}{Portable Document Format}
\nomenclature{PNG}{Portable Network Graphics}
%\nomenclature{XML}{Extensible Markup Language}


% Formelzeichen %%%%%%%%%%%%%%%%%%%%%%%%%%%%%%%%%%%%
\nomenclature[a]{$a$}{Beschleunigung}
\nomenclature[F]{$F$}{Kraft}
\nomenclature[m]{$m$}{Masse}
\nomenclature[P]{$P$}{Leistung}
\nomenclature[U]{$U$}{Spannung}
\nomenclature[R]{$R$}{Widerstand}


				% Datei mit allgemeinen Abkürzungen laden
\renewcommand{\nomname}{Verzeichnis verwendeter Formelzeichen und Abkürzungen}
\setlength{\nomlabelwidth}{.20\hsize}
\renewcommand{\nomlabel}[1]{#1 \dotfill}
\setlength{\nomitemsep}{-\parsep}
\printnomenclature						% Erzeugen des Abkürzungsverzeichnises, siehe auch Inhalt der Datei pages/abkuerzungen.tex
\cleardoublepage
%\renewcommand{\glossaryname}{Glossar}
%\printglossaries
%\cleardoublepage
\listoffigures 							% Erzeugen des Abbildungsverzeichnisses 
\cleardoublepage
\listoftables 							% Erzeugen des Tabellenverzeichnisses
\cleardoublepage

% -----Anhang ------------------------

\appendix
\clearpage
%\pagenumbering{Roman}					% große, römische Seitenzahlen für Anhang, falls gewünscht
%\addchap{Anhang A}
\setcounter{chapter}{1}

\section{Details zu bestimmten theoretischen Grundlagen}

\section{Weitere Details, welche im Hauptteil den Lesefluss behindern}

\addchap{Anhang B}
\setcounter{chapter}{2}
\setcounter{section}{0}
\setcounter{table}{0}
\setcounter{figure}{0}

\section{Versuchsanordnung}

\section{Liste der verwendeten Messgeräte}

\section{Übersicht der Messergebnisse}

\section{Schaltplan und Bild der Prototypenplatine}

\clearpage

Diese Seite wurde eingefügt, um zu zeigen, wie sich der Inhalt der Kopfzeile automatisch füllt.

\addchap{Anhang C}
\setcounter{chapter}{3}
\setcounter{section}{0}
\setcounter{table}{0}
\setcounter{figure}{0}

\section{Struktogramm des Programmentwurfs}

\section{Wichtige Teile des Quellcodes}

\addchap{Anhang D}
\setcounter{chapter}{4}
\setcounter{section}{0}
\setcounter{table}{0}
\setcounter{figure}{0}

\section{Einbinden von PDF-Seiten aus anderen Dokumenten}

Auf den folgenden Seiten wird eine Möglichkeit gezeigt, wie aus einem anderen PDF-Dokument komplette Seiten übernommen werden können. Der Nachteil dieser Methode besteht darin, dass sämtliche Formateinstellungen (Kopfzeilen, Seitenzahlen, Ränder, etc.) auf diesen Seiten nicht angezeigt werden. Die Methode wird deshalb eher selten gewählt. Immerhin sorgt das Package \textit{\glqq pdfpages\grqq}~für eine korrekte Seitenzahleinstellung auf den im Anschluss folgenden \glqq nativen\grqq~\LaTeX-Seiten.

Eine bessere Alternative ist, einzelne Seiten mit \textit{\glqq$\backslash$includegraphics\grqq}~einzubinden. Z.B. wenn Inhalte von Datenblättern wiedergegeben werden sollen.

\includepdf[pages={2-4}]{docs/EingebundenesPDF.pdf}

%\addchap{Abbildungen}
\title{Abbildungen}
\listoffigures		% Zeile auskommentieren bei finalem Dokument!

\end{document}